\section{符号说明}

文中主要符号及含义如表\ref{tab:symbols}所示，
各问题中特有的参数和辅助变量在相应模型中另行说明。

\begin{table}[H]
	\centering
	\caption{主要符号说明}
	\label{tab:symbols}
	
	% 压缩表格字体和行距
	\small
	\renewcommand{\arraystretch}{0.85}
	
	\begin{tabularx}{0.96\textwidth}{
			>{\centering\arraybackslash}p{0.20\textwidth}
			>{\raggedright\arraybackslash}X
			>{\centering\arraybackslash}p{0.15\textwidth}}
		\toprule
		\textbf{符号} & \multicolumn{1}{c}{\textbf{含义}} & \textbf{单位} \\
		\midrule
		
		$i,r,t,s$
		& 任务、区域、时段及候选开工时刻索引
		& -- \\
		
		$a_i,e_i,o_i,d_i$
		& 到达、最早开工、最晚完成时刻及持续时间
		& h \\
		
		$g_i$
		& 任务$i$的GPU需求量
		& GPU \\
		
		$p_i$
		& 任务$i$的AI IT功率
		& MW \\
		
		$l_{ir},\,l_i^{\max}$
		& 执行时延及最大允许时延
		& ms \\
		
		$x_{irs}$
		& 任务$i$是否在区域$r$、时刻$s$开工
		& $\{0,1\}$ \\
		
		$\rho_{ist}$
		& 任务$i$在时段$t$内的运行占比
		& -- \\
		
		$P^{AI}_{rt},\,P^{NonAI}_{rt}$
		& AI IT负荷与非AI IT负荷
		& MW \\
		
		$F_{rt}$
		& 区域$r$的设施侧总负荷
		& MW \\
		
		$G^{ava}_{rt}$
		& 区域$r$的可用GPU容量
		& GPU \\
		
		$PUE_r$
		& 区域$r$的PUE
		& -- \\
		
		$P^{RE}_{rt}$
		& 区域$r$的可用新能源功率
		& MW \\
		
		$\pi_{rt},\,\sigma_{rt}$
		& 购电价格与售电价格
		& 元/MWh \\
		
		$\kappa_{rt}$
		& 电网碳强度
		& tCO$_2$/MWh \\
		
		$B_{rt},\,X_{rt},\,R_{rt}$
		& 电网购电、新能源外送及弃用功率
		& MW \\
		
		$S_{rt}$
		& 时段$t$开始时的储能电量
		& MWh \\
		
		\bottomrule
	\end{tabularx}
\end{table}

